\par\SourceLeaf{425}%
\noindent\hspace*{1.2em}Wir nehmen an, dass die Punkte so geordnet sind, dass sie in demselben Sinn um das von ihnen eingeschlossene Dreieck laufen, wie die Punkte $(1),(2),(3)$; es sei ferner $\lambda$ derjenige Pol des grössten Kreises $L'L''$, welcher auf derselben Seite liegt wie $L$. Man hat dann aus unserm vorigen Lemma
\[
\begin{aligned}
y'z''-z'y''&=\sin L'L''.\cos\lambda(1)\\
z'x''-x'z''&=\sin L'L''.\cos\lambda(2)\\
x'y''-y'x''&=\sin L'L''.\cos\lambda(3).
\end{aligned}
\]
Multiplicirt man also diese Gleichungen mit $x,y,z$ und addirt die Producte, so wird
\[
xy'z''+x'y''z+x''yz'-xy''z'-x'yz''-x''y'z=\sin L'L''.\cos\lambda L,
\]
wofür man nach den bekannten Lehren der sphärischen Trigonometrie auch
\[
\begin{aligned}
&\sin L'L''.\sin LL''.\sin L'\\
&=\sin L'L''.\sin LL'.\sin L''\\
&=\sin L'L''.\sin L'L''.\sin L
\end{aligned}
\]
schreiben kann, wenn $L,L',L''$ die drei Winkel des sphärischen Dreiecks $LL'L''$ bedeuten. Man überzeugt sich zugleich leicht, dass diese Grösse der sechste Theil der Pyramide ist, die zwischen dem Centrum der Kugel und den Punkten $L,L',L''$ gebildet wird (und zwar \textls[100]{positiv}, wenn etc.).

\begin{center}8.\end{center}
\noindent\hspace*{1.2em}Die Natur einer krummen Fläche wird durch eine Gleichung zwischen den Coordinaten ihrer Punkte dargestellt, die wir durch
\[
f(x,y,z)=0
\]
vorstellen; es sei das vollständige Differential von $f(x,y,z)$
\[
=Pdx+Qdy+Rdz,
\]
wo $P,Q,R$ Functionen von $x,y,z$ sein werden. Wir werden immer an der Fläche zwei Seiten unterscheiden, wovon wir die eine die obere, die andere die untere nennen wollen; allgemein zu reden wird beim Durchgang durch
